\section{Conclusiones}
\label{sec:conclusiones}

La evaluación base confirmó que SOS Clinic priorizaba la funcionalidad sobre la
seguridad: las APIs clínicas no autenticaban, las rutas administrativas no
autorizaban por rol y el secreto JWT estaba versionado. El riesgo global medido en
la rama \texttt{pruebas-seguridad-2} era \textbf{crítico}. Después de las
correcciones aplicadas en \texttt{main}, \textbf{22 hallazgos quedaron remediados,
1 mitigado (H-05), 1 parcial (H-18) y 1 documentado (H-25)}. El riesgo residual baja
de forma significativa, pero no desaparece.

Los objetivos de la evaluación se cumplieron: se identificaron secretos y
configuraciones inseguras, se verificó la ausencia de autenticación y autorización
en los servicios de negocio, se demostró el XSS almacenado y el abuso del agente
IA, y se consolidaron 25 temas únicos con evidencia reproducible. La matriz de
hallazgos, las capturas y los resultados crudos permiten repetir las pruebas; la
sección~\ref{sec:estado-remediacion} añade el estado de las correcciones y las
pruebas de verificación ejecutadas en \texttt{main}.

También quedan limitaciones claras. El entorno es local y de desarrollo; ZAP se
ejecutó en modo pasivo; el SCA dependió del caché local de Maven; Trivy no evaluó
\texttt{docker-compose.yml}; y la verificación posterior se apoya en \texttt{mvn
test}, Trivy, gitleaks y pruebas E2E, no en una auditoría independiente. Por tanto,
los resultados no deben interpretarse como una certificación de producción ni como
una garantía de ausencia de vulnerabilidades.

La recomendación final es mantener la seguridad como una propiedad transversal y
cerrar los riesgos residuales: revocación de sesiones o logout con estado, rate
limiting distribuido, autorización por propiedad de cada episodio, CSP estricta sin
\texttt{unsafe-inline}, TLS con HSTS en producción y rotación de los secretos que
permanecen en el historial de Git. Con esas acciones, y una re-evaluación
independiente, el prototipo puede pasar de un riesgo crítico a un riesgo controlado.
